\documentclass[../../main.tex]{subfiles}
\begin{document}

\begin{center}
  {\fontsize{48pt}{58pt}\selectfont \textbf{1997}\normalsize\textbf{年度}}
\end{center}

\begin{table}[htb]
  \centering
  \caption{出題テーマと難易度}
  \begin{tabular}{|c|p{9cm}|c|}
    \hline
    \textbf{問題番号} & \textbf{テーマ} & \textbf{難易度} \\
    \hline
    第1問 & 円上の点から定まる軌跡 & \\
    \hline
    第2問 & 二項係数の最大公約数 & \\
    \hline
    第3問 & 放物線で囲む面積の二等分 & \\
    \hline
    第4問 & 正弦関数の積分と和の不等式 & \\
    \hline
    第5問 & ゲームの得点期待値と勝率 & \\
    \hline
    第6問 & 余弦曲線の接線で囲む面積 & \\
    \hline
  \end{tabular}
\end{table}

\begin{table}[htb]
  \centering
  \caption{解答}
  \renewcommand{\arraystretch}{1.4}
  \begin{tabular}{|c|c|p{9cm}|}
    \hline
    \textbf{問題番号} & \textbf{小問} & \textbf{答え} \\
    \hline
    第1問 & (1) & $t=2$ \\
    \cline{2-3}
    & (2) & $x^2+y^2=4$のうち$x\ge1$の弧（端点を含む） \\
    \hline
    第2問 & -- & 最大公約数は$1$（証明問題） \\
    \hline
    第3問 & (1) & $a^2=\dfrac{k-\left(\dfrac{k^4}{4(k-1)}\right)^{1/3}}{k-1}$ \\
    \cline{2-3}
    & (2) & $1-4^{-1/3}$ \\
    \hline
    第4問 & (1) & 指定された積分不等式が成立（証明問題） \\
    \cline{2-3}
    & (2) & $\sum_{k=1}^7\sin\dfrac{k\pi}{8}<\dfrac{16}{\pi}$（証明問題） \\
    \hline
    第5問 & (1) & $E_A=-2p^3+3p^2-4p+6$，$E_B=-12p^3+21p^2-12p+6$．$E_A>E_B\iff p<4/5$ \\
    \cline{2-3}
    & (2) & 言えない．$P_A>P_B\iff p<2-\sqrt{2}$ \\
    \hline
    第6問 & (1) & $S(t)=\dfrac{(\cos t+t\sin t)^2}{2\sin t}$ \\
    \cline{2-3}
    & (2) & $\pi/4<t_0<1$で唯一の最小値 \\
    \cline{2-3}
    & (3) & $S(t_0)=2t_0\cos t_0>\dfrac{\sqrt{2}\pi}{4}$ \\
    \hline
  \end{tabular}
\end{table}

\paragraph{試験形態}
6問．試験時間は要確認．

\end{document}
